\section{Variance, standard deviation, and moments}

Expectation describes the center of a probability distribution, but it does not
describe how spread out the values are. Two random variables may have the same
expectation and very different variability. Variance is the basic measure of
spread around the expectation.

Let
\[
\mu=\E[X].
\]
The quantity
\[
X-\mu
\]
measures the deviation of \(X\) from its mean. Positive values are above the
mean; negative values are below the mean.

\subsection*{Why square the deviation?}

One might first try to average the deviations:
\[
\E[X-\mu].
\]
But by linearity of expectation,
\[
\E[X-\mu]=\E[X]-\mu=0.
\]
The positive and negative deviations cancel. This does not measure spread.

To avoid cancellation, we square the deviations. Large deviations then contribute
more strongly, and all contributions are non-negative.

\begin{definitionbox}
The \textbf{variance} of a random variable \(X\) is
\[
\Var(X)=\E[(X-\E[X])^2],
\]
provided this expectation exists.
\end{definitionbox}

For a discrete random variable, this becomes
\[
\Var(X)=\sum_{x\in\operatorname{supp}(X)}(x-\mu)^2p_X(x),
\]
where \(\mu=\E[X]\).

\subsection*{Standard deviation}

Variance is measured in squared units. If \(X\) is measured in euros, then
\(\Var(X)\) is measured in squared euros. To return to the original units, we take
the square root.

\begin{definitionbox}
The \textbf{standard deviation} of \(X\) is
\[
\sigma_X=\sqrt{\Var(X)}.
\]
\end{definitionbox}

The standard deviation is a typical scale of deviation from the mean. It is not
the same as the average absolute distance, but it is one of the most important
measures of spread in probability and statistics.

\subsection*{A useful computational formula}

The definition
\[
\Var(X)=\E[(X-\mu)^2]
\]
is conceptually clear. For calculations, another form is often easier. Expand
the square:
\[
(X-\mu)^2=X^2-2\mu X+\mu^2.
\]
Taking expectation gives
\[
\E[(X-\mu)^2]=\E[X^2-2\mu X+\mu^2].
\]
Using linearity,
\[
\E[(X-\mu)^2]=\E[X^2]-2\mu\E[X]+\mu^2.
\]
Since \(\mu=\E[X]\), we get
\[
\E[(X-\mu)^2]=\E[X^2]-2\mu^2+\mu^2.
\]
Therefore
\[
\Var(X)=\E[X^2]-(\E[X])^2.
\]

\begin{theorembox}
For every random variable for which the quantities are defined,
\[
\Var(X)=\E[X^2]-(\E[X])^2.
\]
\end{theorembox}

This formula should not replace the meaning of variance. It is a computational
consequence of the definition.

\subsection*{Example: one fair die}

Let \(X\) be the result of one fair die roll. We already know
\[
\E[X]=\frac{1+2+3+4+5+6}{6}=\frac72.
\]
Now compute
\[
\E[X^2]=\frac{1^2+2^2+3^2+4^2+5^2+6^2}{6}.
\]
The numerator is
\[
1+4+9+16+25+36=91.
\]
Thus
\[
\E[X^2]=\frac{91}{6}.
\]
Therefore
\[
\Var(X)=\frac{91}{6}-\left(\frac72\right)^2
=\frac{91}{6}-\frac{49}{4}.
\]
Using denominator \(12\),
\[
\Var(X)=\frac{182}{12}-\frac{147}{12}=\frac{35}{12}.
\]
The standard deviation is
\[
\sigma_X=\sqrt{\frac{35}{12}}.
\]

\subsection*{Example: an indicator}

Let \(I_A\) be the indicator of an event \(A\), and let
\[
p=P(A).
\]
Then
\[
\E[I_A]=p.
\]
Because \(I_A\) takes only the values \(0\) and \(1\), we have
\[
I_A^2=I_A.
\]
Hence
\[
\E[I_A^2]=\E[I_A]=p.
\]
Therefore
\[
\Var(I_A)=p-p^2=p(1-p).
\]

This result is important because many counts are sums of indicators. It also
shows that the variance of a zero-one variable is largest when \(p=1/2\), and
small when the event is almost impossible or almost certain.

\subsection*{Moments}

Expectations of powers of \(X\) are called moments. The \(k\)-th raw moment is
\[
\E[X^k].
\]
For a discrete random variable,
\[
\E[X^k]=\sum_x x^k p_X(x),
\]
provided the sum is defined.

The first raw moment is the expectation:
\[
\E[X].
\]
The second raw moment \(\E[X^2]\) is used in the computation of variance. Higher
moments describe further aspects of the shape of the distribution, but in this
course the most important quantities at this stage are expectation and variance.

\subsection*{Interpreting variance carefully}

Variance is always non-negative because it is the expectation of a square:
\[
\Var(X)=\E[(X-\mu)^2]\geq 0.
\]
It is zero exactly when \(X\) is constant with probability one. If there is no
spread, every possible value is equal to the mean.

A larger variance indicates more spread, but it must be interpreted relative to
the scale of the problem. A variance of \(100\) is large for exam scores out of
\(10\), but small for yearly income measured in euros.

\begin{summarybox}
Expectation measures center. Variance measures average squared distance from the
center. Standard deviation returns this spread to the original units:
\[
\sigma_X=\sqrt{\Var(X)}.
\]
Together, expectation and variance give the first basic numerical summary of a
probability law.
\end{summarybox}
